\documentclass[10pt]{article}
\usepackage[margin=2.2cm]{geometry}
\usepackage[T1]{fontenc}
\usepackage{lmodern}
\usepackage{enumitem}
\usepackage{hyperref}
\hypersetup{pdfauthor={Anonymous}, pdftitle={Summary of changes and response to the prior review}}
\pdfinfoomitdate=1
\pdfsuppressptexinfo=-1
\pdftrailerid{}
\setlist{nosep,leftmargin=1.4em}
\setlength{\parskip}{0.45em}
\setlength{\parindent}{0pt}

\title{Summary of Changes and Response to the Prior Review\\[0.3em]
\large Submission: ``What Prompt Mediation Stops in LLM Agent Pipelines, and How to Measure It''}
\author{Anonymous submission}
\date{}

\begin{document}
\maketitle

\section*{Prior review}
An earlier version of this work was submitted to ACM Transactions on Privacy and Security (submission
TOPS-2026-06-0395) under a different title and was rejected on the Associate Editor's recommendation. As the
PoPETs submission guidelines require, the decision letter and the Associate Editor's comments are attached
unchanged, and this document summarizes what changed and answers each comment. Section numbers refer to the
present submission; appendix sections are named by title.

\section*{Overview of the changes}
The paper has been rebuilt rather than revised. It no longer presents a new privacy mechanism. It is now a
measurement study of what ingress prompt mediation can and cannot stop in LLM agent pipelines, together with a
validation of the instruments used to score leakage. In particular:
\begin{itemize}
\item \textbf{Framing and title.} The contribution is the measurement (\S1): identifier exposure across agent
  boundaries under one protocol against installed baselines; where mediation is applied in a live tool-using
  agent; factual leakage measured with an adversary and with blind human labels; and the agreement of automated
  leakage measures with people. Replacement with restoration is stated as prior art (\S3).
\item \textbf{The constructed benchmark is gone.} All results built on the 256-prompt benchmark of the earlier
  version, and its figures, have been removed. Every reported result now uses public benchmarks (AI4Privacy,
  PrivacyLens, AgentDojo) or public real-text transfer slices, and \S5.1 carries the extrapolation paragraph the
  call requires for synthetic evaluations.
\item \textbf{Real agents.} Two open-weight agents produce the downstream actions (\S6.3); a live tool-using
  stack, in which the agent itself writes every boundary artifact, varies where mediation is applied (\S6.5); and
  AgentDojo tools execute against their sandboxed environments under prompt injection (\S6.6).
\item \textbf{Strong baselines.} Installed Presidio, a Casper-style sanitizer with its restoration map, a one-shot
  LLM sanitizer, and PrivacyLens's own privacy-enhanced agent prompt, all run under the same protocol.
\item \textbf{Realistic attacks.} Adaptive surface-form evasion against the deployed detectors, including a
  language-model attacker with detector feedback (\S6.6); an attribute-recovery attack on the released surface at
  three attacker sizes (\S6.2); prompt injection against a live agent (\S6.6); and an analysis of what the
  replacement tokens themselves disclose (\S4).
\item \textbf{Utility and statistics.} The custom utility rubric is gone. Utility is each benchmark's own score,
  and the benchmark's leakage judge is validated against two annotators. Paired tests use Holm correction,
  intervals resample the evaluation unit, and run-to-run variation is measured (appendix ``Portability to a
  Second Downstream Agent'').
\item \textbf{Reproducibility.} Every table and figure is generated from released per-unit records; the model
  roster gives every model's quantization and digest; per-run manifests record decoding settings, seeds, library
  versions, server settings and hardware.
\end{itemize}

\section*{Response to the Associate Editor's comments}

\paragraph{1. Prior work (PrivacyRestore, Casper) and novelty.}
We agree, and the paper no longer claims a new mechanism. \S3 now covers the closest prompt sanitizers with
restoration (Hide-and-Seek, Casper, PrivacyRestore, PromptGraph, PAAC, Pr$\varepsilon\varepsilon$mpt), model-based
rewriting (PAPILLON, Rescriber, LLM anonymizers), and information-flow control and agent data-flow defenses
(Denning, JIF, TaintDroid, CaMeL, AirGapAgent), and states that replacement with restoration, and deferring
restoration in an agent loop, are prior art. Casper's mechanism is reimplemented and compared under our protocol,
including its client-held restoration map (Table~5); PrivacyRestore is discussed but not run, because it steers
activations inside the serving model, whereas the layer studied here sits in front of a black-box model. The
contribution is the measurement listed above.

\paragraph{2. The 256-prompt benchmark from 32 templates.}
Removed. No result of the paper rests on it. The learned detector is trained on the AI4Privacy training split and
evaluated on 300 held-out documents; that it is evaluated in domain there, and out of domain on PrivacyLens, is
stated in \S5.3 and the Limitations, together with the fact that its margin over the Casper-style detector
disappears out of domain.

\paragraph{3. The deterministic reference agent, and contextual disclosure.}
Both concerns are now addressed by measurements rather than argument. Agent behavior is measured where a model
writes the artifact: identifiers in the final action of two agents (\S6.3), and a live tool-using agent with a
retrieval index, a note store and an application log, under six mediation placements and two detector arms
(\S6.5). There, mediating only the prompt leaves 61--68\% of the values that tools return in the agent's context;
only mediating every read edge keeps them out. Contextual disclosure is measured directly (\S6.2): an adversary
choosing among four same-kind candidates recovers 64\% of the values the strongest detector removed (chance 25\%),
at every attacker size tested, and blind human labels on 240 items, checked by a second reader at
$\kappa=0.849$, show what remains inferable. The paper states that SPE measures literal survival only, and its
central result is precisely the gap between literal and factual leakage.

\paragraph{4. Enterprise redaction, AC/TSR, and the custom utility rubric.}
The comparison and the rubric are removed. Utility is now each benchmark's own instrument: PrivacyLens's
helpfulness score from its original judge and verbatim prompts, and AgentDojo's state-checked task success. The
benchmark's leakage judge is validated against two annotators (appendix ``Validating the Benchmark's Own
Judge''). Utility costs are reported where they are large, including a helpfulness drop from 2.46 to 1.40 for
the learned detector on the 70B agent, and differences below run-to-run variation are reported as unresolved.

\paragraph{5. Detector precision, model identities, OCR and runtime records.}
Precision is now a column of the exposure table (Table~3): the learned detector's span precision is 0.35,
against 0.66--0.78 for the transformer, Casper-style and Presidio detectors, and the retention column
prices its over-redaction. The model roster names every model with its parameter count, quantization and digest,
and per-run manifests record decoding parameters, seeds, library versions, server settings and hardware. The OCR
transfer slices report their provenance, and the released records regenerate every reported number.

\paragraph{6. The span-level LDP analysis.}
Removed. \S2 now states that the evaluated pipeline is deterministic and provides no end-to-end formal guarantee,
and that all reported quantities are empirical measurements with explicit scope.

\paragraph{Closing recommendation: real agent stacks, strong baselines, realistic attacks.}
See the overview: the live-stack placement experiment, AgentDojo with executing tools, two downstream agents,
installed and reimplemented baselines including a prompt-only defense, and adaptive, attribute-recovery and
injection attacks.

\section*{Other substantive changes}
\begin{itemize}
\item Trust domains, the trusted computing base and the authorized boundary are defined (\S2), and the
  architecture figure is redrawn programmatically to match the text.
\item Metric definitions match the code that computes them: PER is a pooled ratio with effective counts, attack
  success is counted per span, fact retention matches whole words and is described as a permissive lexical
  measure.
\item The restoration-timing contrast is described for what it is: deferral against restoring before planning.
\item Every bibliographic entry was checked against its canonical record.
\item The ethics, open-science and AI-use sections follow the PoPETs 2027 guidelines.
\end{itemize}

\end{document}
